\documentclass[12pt,a4paper]{article}

\usepackage[utf8]{inputenc}
\usepackage[T1]{fontenc}
\usepackage{times}
\usepackage{setspace}
\usepackage{indentfirst}
\usepackage{geometry}
\usepackage{listings}
\usepackage{xcolor}
\usepackage{hyperref}

\geometry{
a4paper,
top=3cm,
bottom=2cm,
left=3cm,
right=2cm
}

\onehalfspacing
\setlength{\parindent}{1.25cm}

\lstset{
basicstyle=\ttfamily\small,
breaklines=true,
frame=single,
numbers=left,
numberstyle=\tiny,
showstringspaces=false,
columns=fullflexible
}

\begin{document}

\begin{titlepage}
\centering

```
{\large UNINASSAU\par}
\vspace{0.5cm}
{\large CURSO DE CIÊNCIA DA COMPUTAÇÃO\par}
\vspace{0.5cm}
{\large DISCIPLINA: COMPILADORES\par}

\vfill

{\bfseries\Large RELATÓRIO TÉCNICO\par}
\vspace{0.5cm}
{\bfseries\Large COMPILADOR E INTERPRETADOR DA LINGUAGEM MMRLANG\par}

\vfill

{\large Integrantes:\par}
\vspace{0.3cm}

{\small Adriano Mikhael M L Cavalcante -- 01605455 -- \url{https://github.com/AdrianoMikhael}\par}
{\small João Lucas Vital -- 01624753 -- \url{https://github.com/JoaoLucasVital}\par}
{\small Lucas Soares Silva -- 01631745 -- \url{https://github.com/lqnhs1}\par}

\vfill

{\large Maceió\par}
{\large 2026\par}
```

\end{titlepage}

\begin{titlepage}
\centering

```
{\large Adriano Mikhael M L Cavalcante -- 01605455\par}
{\large João Lucas Vital -- 01624753\par}
{\large Lucas Soares Silva -- 01631745\par}

\vfill

{\bfseries\Large COMPILADOR E INTERPRETADOR DA LINGUAGEM MMRLANG\par}

\vspace{3cm}

\begin{flushright}
\begin{minipage}{8cm}
Relatório técnico apresentado à disciplina de Compiladores como requisito avaliativo da AV2, com o objetivo de demonstrar o desenvolvimento de um compilador e interpretador simples, utilizando versionamento no GitHub e documentação em LaTeX.
\end{minipage}
\end{flushright}

\vfill

{\large Maceió\par}
{\large 2026\par}
```

\end{titlepage}

\section{Resumo}

Este relatório apresenta o desenvolvimento da MMRLang, uma linguagem de programação educacional criada para demonstrar conceitos fundamentais da disciplina de Compiladores.

O projeto consiste em um compilador e interpretador simples, desenvolvido em Python, capaz de ler arquivos com extensão .mmr, interpretar comandos próprios da linguagem, armazenar variáveis, avaliar expressões, executar estruturas condicionais, realizar repetições e apresentar mensagens de erro.

Além do desenvolvimento do interpretador, o projeto também utiliza documentação em LaTeX, edição por meio da IDE Texmaker, compilação com MiKTeX e versionamento do código no GitHub com colaboração entre os integrantes.

\section{Introdução}

A disciplina de Compiladores estuda os processos utilizados para analisar, interpretar ou transformar códigos escritos em uma linguagem de programação. Esses processos envolvem etapas como leitura do código-fonte, identificação de comandos, avaliação de expressões, execução de instruções e tratamento de erros.

Neste trabalho, foi desenvolvida a MMRLang, uma linguagem criada pela equipe com comandos próprios e sintaxe em português. A linguagem foi pensada para ser simples, didática e levemente descontraída, mas mantendo conceitos importantes presentes em linguagens reais.

O interpretador foi desenvolvido em Python, por ser uma linguagem de fácil leitura e adequada para projetos acadêmicos. O relatório foi produzido em LaTeX, utilizando o MiKTeX como compilador tipográfico e o Texmaker como ambiente de edição.

\section{Objetivos}

\subsection{Objetivo Geral}

Desenvolver um compilador e interpretador simples para uma linguagem própria, aplicando conceitos estudados na disciplina de Compiladores e documentando o projeto em LaTeX.

\subsection{Objetivos Específicos}

\begin{itemize}
\item Criar uma linguagem de programação educacional chamada MMRLang;
\item Desenvolver um interpretador em Python;
\item Implementar variáveis, saída de dados, condições e repetições;
\item Permitir a execução de arquivos com extensão .mmr;
\item Tratar erros comuns encontrados no código-fonte;
\item Organizar programas de teste em arquivos separados;
\item Versionar o projeto no GitHub;
\item Produzir relatório técnico em LaTeX.
\end{itemize}

\section{Significado da MMRLang}

O nome MMRLang é formado por duas partes: MMR e Lang.

A sigla MMR representa a identidade escolhida pela equipe para a linguagem. Já o termo Lang é uma abreviação da palavra inglesa language, que significa linguagem.

Portanto, MMRLang pode ser entendida como Linguagem MMR.

A proposta da linguagem é unir uma estrutura lógica de programação com comandos simples, em português e com uma identidade própria.

\section{Principais Características}

A MMRLang possui as seguintes características:

\begin{itemize}
\item Sintaxe simples e em português;
\item Arquivos com extensão .mmr;
\item Interpretador desenvolvido em Python;
\item Comandos próprios e fáceis de entender;
\item Suporte a variáveis;
\item Suporte a textos, números e valores lógicos;
\item Suporte a operações matemáticas;
\item Suporte a estruturas condicionais;
\item Suporte a estruturas de repetição;
\item Tratamento de erros com mensagens claras;
\item Testes separados por arquivos.
\end{itemize}

\section{Tecnologias Utilizadas}

\subsection{Python}

Python foi utilizado para desenvolver o interpretador da MMRLang. A escolha ocorreu por sua sintaxe clara, facilidade de leitura de arquivos e boa aplicação em projetos acadêmicos.

\subsection{LaTeX}

LaTeX foi utilizado para escrever o relatório técnico do projeto. Ele permite criar documentos bem estruturados, com seções, sumário e blocos de código.

\subsection{MiKTeX}

O MiKTeX foi utilizado como compilador LaTeX. Ele é responsável por transformar o arquivo main.tex em um documento PDF.

\subsection{Texmaker}

O Texmaker foi utilizado como IDE para editar e compilar o relatório em LaTeX.

\subsection{GitHub}

O GitHub foi utilizado para versionamento do projeto e colaboração entre os integrantes da equipe.

\section{Estrutura do Projeto}

A estrutura do projeto foi organizada da seguinte forma:

\begin{lstlisting}
compilador-av2/
|-- main.py
|-- testes/
|   |-- programa1.mmr
|   |-- programa2.mmr
|   |-- programa3.mmr
|   |-- programa4.mmr
|   |-- programa5.mmr
|-- relatorio/
|   |-- main.tex
|-- README.md
|-- .gitignore
\end{lstlisting}

\section{Comandos da MMRLang}

A MMRLang usa comandos próprios para representar ações comuns em programação. Mesmo com nomes diferentes, esses comandos possuem funções semelhantes às de linguagens reais.

\subsection{Lista de Comandos}

\begin{itemize}
\item guarda: cria ou altera uma variável. Equivalente a var, let ou atribuição.
\item falaai: exibe uma informação na tela. Equivalente a print ou console.log.
\item eai: inicia uma condição. Equivalente ao if.
\item senaobicho: executa um bloco caso a condição seja falsa. Equivalente ao else.
\item fimtrem: finaliza um bloco condicional.
\item repete: repete um bloco de comandos. Equivalente a um laço de repetição.
\item acaboupo: finaliza um bloco de repetição.
\item verdadeiro: representa o valor lógico verdadeiro.
\item falso: representa o valor lógico falso.
\item e: representa o operador lógico E.
\item ou: representa o operador lógico OU.
\item nao: representa o operador lógico de negação.
\end{itemize}

\section{Explicação dos Nomes dos Comandos}

\subsection{guarda}

O comando guarda foi escolhido porque sua função é guardar um valor dentro de uma variável.

\begin{lstlisting}
guarda nome = "Maria";
guarda idade = 20;
\end{lstlisting}

Nesse exemplo, a MMRLang guarda o texto Maria na variável nome e o número 20 na variável idade.

\subsection{falaai}

O comando falaai representa a saída de dados. Ele foi escolhido por ser uma forma simples e descontraída de indicar que o programa deve mostrar algo na tela.

\begin{lstlisting}
falaai "Ola, mundo!";
falaai nome;
\end{lstlisting}

Esse comando funciona de maneira semelhante ao print em Python.

\subsection{eai}

O comando eai representa uma condição. Ele faz uma pergunta lógica ao programa.

\begin{lstlisting}
eai idade >= 18
falaai "Maior de idade";
senaobicho
falaai "Menor de idade";
fimtrem
\end{lstlisting}

A ideia é semelhante a perguntar se aquela condição é verdadeira. Se for verdadeira, o primeiro bloco é executado. Caso contrário, o bloco do senaobicho é executado.

\subsection{senaobicho}

O comando senaobicho representa o caso contrário de uma condição. Ele funciona como o else em outras linguagens de programação.

\begin{lstlisting}
eai nota >= 7
falaai "Aprovado";
senaobicho
falaai "Recuperacao";
fimtrem
\end{lstlisting}

\subsection{fimtrem}

O comando fimtrem finaliza a estrutura condicional iniciada com eai. Ele indica ao interpretador que o bloco de condição terminou.

\begin{lstlisting}
eai idade >= 18
falaai "Autorizado";
senaobicho
falaai "Negado";
fimtrem
\end{lstlisting}

\subsection{repete}

O comando repete executa um bloco de código várias vezes.

\begin{lstlisting}
repete 3
falaai "MMRLang em execucao";
acaboupo
\end{lstlisting}

Nesse caso, a mensagem será exibida três vezes.

\subsection{acaboupo}

O comando acaboupo finaliza um bloco de repetição iniciado com repete.

\begin{lstlisting}
repete 5
falaai "Repetindo";
acaboupo
\end{lstlisting}

\section{Operadores Suportados}

A MMRLang suporta operadores matemáticos, operadores de comparação e operadores lógicos.

\subsection{Operadores Matemáticos}

\begin{itemize}
\item Soma: sinal de mais.
\item Subtração: sinal de menos.
\item Multiplicação: asterisco.
\item Divisão: barra.
\item Resto da divisão: porcentagem.
\end{itemize}

\subsection{Operadores de Comparação}

\begin{itemize}
\item Igual.
\item Diferente.
\item Maior.
\item Menor.
\item Maior ou igual.
\item Menor ou igual.
\end{itemize}

\subsection{Operadores Lógicos}

\begin{itemize}
\item e: operador lógico E.
\item ou: operador lógico OU.
\item nao: operador lógico de negação.
\end{itemize}

\section{Funcionamento do Interpretador}

O interpretador da MMRLang executa as seguintes etapas:

\begin{enumerate}
\item Recebe o caminho de um arquivo .mmr;
\item Verifica se o arquivo existe;
\item Verifica se a extensão do arquivo é válida;
\item Lê o conteúdo do arquivo linha por linha;
\item Remove comentários, espaços desnecessários e ponto e vírgula final;
\item Identifica os comandos da linguagem;
\item Executa variáveis, saídas, condições e repetições;
\item Avalia expressões matemáticas e lógicas;
\item Exibe mensagens de erro caso encontre problemas.
\end{enumerate}

\section{Exemplo de Código em MMRLang}

\begin{lstlisting}
guarda nome = "Maria";
guarda idade = 20;
guarda nota = 8;

falaai "Sistema MMR iniciado";
falaai "Nome do aluno:";
falaai nome;

eai idade >= 18
falaai "Usuario maior de idade";
senaobicho
falaai "Usuario menor de idade";
fimtrem

eai nota >= 7
falaai "Aluno aprovado";
senaobicho
falaai "Aluno em recuperacao";
fimtrem

repete 3
falaai "Finalizando processamento...";
acaboupo

falaai "Programa encerrado";
\end{lstlisting}

\section{Saída Esperada}

\begin{lstlisting}
Sistema MMR iniciado
Nome do aluno:
Maria
Usuario maior de idade
Aluno aprovado
Finalizando processamento...
Finalizando processamento...
Finalizando processamento...
Programa encerrado
\end{lstlisting}

\section{Programas de Teste}

Os testes foram separados em arquivos diferentes dentro da pasta testes. Essa organização facilita a execução individual de cada funcionalidade da linguagem.

\begin{itemize}
\item programa1.mmr: testa variáveis e saída de dados.
\item programa2.mmr: testa operações matemáticas.
\item programa3.mmr: testa estrutura condicional.
\item programa4.mmr: testa estrutura de repetição.
\item programa5.mmr: testa um programa completo com variáveis, condições e repetição.
\end{itemize}

\section{Tratamento de Erros}

O tratamento de erros é uma parte importante do interpretador da MMRLang. Ele permite que o programa identifique problemas no código-fonte e mostre mensagens mais claras para o usuário.

\subsection{Erros Apresentados}

\begin{itemize}
\item Arquivo não encontrado: acontece quando o arquivo informado pelo usuário não existe.
\item Extensão inválida: acontece quando o arquivo não possui extensão .mmr.
\item Comando desconhecido: acontece quando é usado um comando que não existe na MMRLang.
\item Variável não declarada: acontece quando uma variável é usada antes de ser criada.
\item Nome de variável inválido: acontece quando o nome da variável não segue o formato permitido.
\item Expressão vazia: acontece quando um comando precisa de valor, mas nenhum valor é informado.
\item Erro no comando guarda: acontece quando a atribuição não utiliza o sinal de igual.
\item Erro no comando eai: acontece quando a condição não é informada.
\item Bloco eai não finalizado: acontece quando falta o comando fimtrem.
\item Bloco repete não finalizado: acontece quando falta o comando acaboupo.
\item Quantidade inválida no repete: acontece quando o valor da repetição não é um número inteiro.
\item Repetição negativa: acontece quando a quantidade de repetições é menor que zero.
\item Divisão por zero: acontece quando uma expressão tenta dividir um número por zero.
\item Operador não suportado: acontece quando é usado um operador não permitido pela linguagem.
\item Fechamento fora de bloco: acontece quando fimtrem, senaobicho ou acaboupo aparecem fora do local correto.
\end{itemize}

\subsection{Exemplos de Erros}

\subsubsection{Variável não declarada}

\begin{lstlisting}
falaai nome;
\end{lstlisting}

Mensagem esperada:

\begin{lstlisting}
Erro MMR:
[Linha 1] A variavel 'nome' nao foi declarada.
\end{lstlisting}

\subsubsection{Comando desconhecido}

\begin{lstlisting}
escreva "Ola";
\end{lstlisting}

Mensagem esperada:

\begin{lstlisting}
Erro MMR:
[Linha 1] Comando desconhecido: escreva "Ola"
\end{lstlisting}

\subsubsection{Bloco condicional não finalizado}

\begin{lstlisting}
eai idade >= 18
falaai "Maior de idade";
\end{lstlisting}

Mensagem esperada:

\begin{lstlisting}
Erro MMR:
[Linha 1] Bloco 'eai' nao foi finalizado com 'fimtrem'.
\end{lstlisting}

\subsubsection{Bloco de repetição não finalizado}

\begin{lstlisting}
repete 3
falaai "Executando";
\end{lstlisting}

Mensagem esperada:

\begin{lstlisting}
Erro MMR:
[Linha 1] Bloco 'repete' nao foi finalizado com 'acaboupo'.
\end{lstlisting}

\section{Como Rodar o Projeto}

Para executar o projeto, é necessário ter o Python instalado no computador.

Primeiro, clone o repositório:

\begin{lstlisting}
git clone LINK_DO_REPOSITORIO
\end{lstlisting}

Depois, entre na pasta do projeto:

\begin{lstlisting}
cd compilador-av2
\end{lstlisting}

Execute o interpretador:

\begin{lstlisting}
python main.py
\end{lstlisting}

Após executar o interpretador, será solicitado o caminho do arquivo:

\begin{lstlisting}
Arquivo MMR:
\end{lstlisting}

Exemplo de arquivo para execução:

\begin{lstlisting}
testes/programa1.mmr
\end{lstlisting}

\section{Relatório em LaTeX}

O relatório do projeto foi produzido em LaTeX, utilizando o compilador MiKTeX e a IDE Texmaker.

O MiKTeX é responsável por compilar o arquivo main.tex e gerar o documento em PDF. Já o Texmaker foi utilizado como editor para escrever, organizar e compilar o relatório.

\section{Versionamento no GitHub}

O projeto foi versionado no GitHub para permitir organização, histórico de alterações e colaboração entre os integrantes.

Todos os integrantes devem ser adicionados como colaboradores no repositório. Dessa forma, cada integrante pode contribuir com alterações no código, testes e documentação.

\subsection{Exemplos de Commits}

\begin{lstlisting}
feat: cria interpretador inicial da MMRLang
feat: adiciona comandos guarda e falaai
feat: implementa estrutura condicional eai
feat: implementa estrutura de repeticao repete
docs: adiciona README do projeto
docs: adiciona relatorio em LaTeX
fix: corrige tratamento de erros
test: adiciona programas de teste
\end{lstlisting}

\section{Uso de Inteligência Artificial}

A inteligência artificial foi utilizada como ferramenta de apoio durante o desenvolvimento do projeto. Ela auxiliou na organização da estrutura do interpretador, na criação da sintaxe da linguagem, na elaboração dos exemplos de teste, na melhoria do tratamento de erros e na escrita da documentação.

As sugestões fornecidas pela IA foram analisadas, testadas e adaptadas pela equipe, garantindo que o projeto atendesse aos requisitos da disciplina e fosse compreendido pelos integrantes.

\section{Dificuldades Encontradas}

Durante o desenvolvimento do projeto, algumas dificuldades foram encontradas, como:

\begin{itemize}
\item Definir uma sintaxe própria para a linguagem;
\item Organizar os blocos condicionais e de repetição;
\item Criar mensagens de erro compreensíveis;
\item Separar os testes por funcionalidade;
\item Organizar o versionamento do código no GitHub;
\item Adaptar o relatório ao formato LaTeX.
\end{itemize}

Essas dificuldades foram resolvidas por meio de testes, revisão do código e organização das responsabilidades entre os integrantes.

\section{Integrantes}

\begin{itemize}
\item Adriano Mikhael M L Cavalcante -- 01605455 -- \url{https://github.com/AdrianoMikhael}
\item João Lucas Vital -- 01624753 -- \url{https://github.com/JoaoLucasVital}
\item Lucas Soares Silva -- 01631745 -- \url{https://github.com/lqnhs1}
\end{itemize}

\section{Conclusão}

O desenvolvimento da MMRLang permitiu aplicar, na prática, conceitos importantes da disciplina de Compiladores. O projeto demonstrou como uma linguagem pode ser lida, interpretada e executada a partir de comandos definidos pela própria equipe.

Apesar de simples, a MMRLang possui características fundamentais de uma linguagem de programação, como variáveis, expressões matemáticas, estruturas condicionais, laços de repetição e tratamento de erros.

Além disso, o uso de LaTeX para a documentação e do GitHub para versionamento contribuiu para uma organização mais profissional do projeto, permitindo colaboração entre os integrantes e registro das etapas de desenvolvimento.

\section{Referências}

\begin{itemize}
\item AHO, Alfred V.; LAM, Monica S.; SETHI, Ravi; ULLMAN, Jeffrey D. Compiladores: princípios, técnicas e ferramentas. Pearson, 2008.
\item Python Software Foundation. Python Documentation. Disponível em: https://docs.python.org/3/.
\item GitHub Docs. Disponível em: https://docs.github.com/.
\item LaTeX Project. LaTeX Documentation. Disponível em: https://www.latex-project.org/help/documentation/.
\item MiKTeX Documentation. Disponível em: https://miktex.org/docs.
\item Texmaker Documentation. Disponível em: https://www.xm1math.net/texmaker/.
\end{itemize}

\end{document}
